% Part: incompleteness
% Chapter: representability-in-q
% Section: beta-function

\documentclass[../../../include/open-logic-section]{subfiles}

\begin{document}

\olfileid{inc}{req}{bet}
\olsection{Lema Fungsi Beta}


Kanggo nuduhake yen kita bisa nindakake rekursi primitif nalika
panjumlahan, perkalian, lan $\Char{=}$ kasedhiya, kita kudu
ngembangake fungsi kang nggarap urutan. (Yen pangkat uga
kasedhiya, tugas iki luwih gampang.) Nalika kita nduweni rekursi
primitif, kita bisa netepake perkara kaya ``prima kaping-$n$,'' lan
milih pangodean kang cukup langsung. Nanging ing kene kita ora
nduweni rekursi primitif---malah kita arep nuduhake yen rekursi
primitif bisa ditindakake nganggo panggolekan tanpa wates---mula kita kudu
nganggo cara kang luwih cerdik.

\begin{lem}
\ollabel{lem:beta}
Ana fungsi $\beta(d,i)$ saengga kanggo saben urutan $a_0$,
\dots,~$a_n$ ana wilangan~$d$ kang ndadekake, kanggo saben $i \le n$,
$\beta(d,i) = a_i$. Kajaba iku, $\beta$ bisa ditegesi saka fungsi
dhasar mung nganggo komposisi lan panggolekan tanpa wates reguler.
\end{lem}

Anggepen $d$ minangka kode urutan $\tuple{a_0, \dots, a_n}$, lan
$\beta(d,i)$ mbalekake unsur kaping-$i$. (Gatekna yen ``pangodean'' iki
\emph{ora} nggunakake pangodean pangkat prima kang wis kita kenal!)
Lema iki winates banget; ora nyatakake yen kita bisa nyambungake
urutan utawa nambah unsur ing pungkasane, utawa malah bisa
\emph{ngitung}~$d$ saka $a_0$, \dots,~$a_n$ nganggo fungsi kang bisa
ditegesi liwat komposisi lan panggolekan tanpa wates reguler. Lema
iki mung nyatakake yen ana fungsi ``dekode'' saengga saben urutan
nduweni ``kode.''

Notasi $\beta$ iki saka G\"odel. Maneh, bab kang angel nalika mbuktekake
lema iki yaiku netepake~$\beta$ kang cocog nganggo sarana kang katone
winates, yaiku mung komposisi lan panggolekan tanpa wates---nanging kita oleh
nggunakake panjumlahan, perkalian, lan~$\Char{=}$. Ana maneka cara
mbuktekake lema iki, nanging salah siji kang paling cetha isih cara
asline G\"odel, kang nggunakake fakta teori wilangan aran Teorema
Sunzi (lumrahe diarani ``Teorema Sisa Cina'').

\begin{defn}
Rong wilangan natural $a$ lan $b$ \emph{relatif prima} yen lan mung yen
faktor bebarengan kang paling gedhe yaiku~$1$; kanthi tembung liya,
ora ana faktor bebarengan liyane.
\end{defn}

\begin{defn}
Wilangan natural $a$ lan $b$ diarani \emph{kongruen modulo~$c$},
$a \equiv b \mod c$, kanthi $c>0$, yen lan mung yen $c \mid (a-b)$, yaiku $a$ lan $b$
nduweni sisa kang padha nalika dipara karo~$c$.
\end{defn}

Iki Teorema Sunzi:
\begin{thm}
Upamana $x_0$, \dots,~$x_n$ positif lan relatif prima saben pasangan. Tetepna
$y_0$, \dots,~$y_n$ minangka wilangan sembarang. Mula ana wilangan $z$ saengga
\begin{align*}
z & \equiv y_0 \mod x_0 \\
z & \equiv y_1 \mod x_1 \\
& \vdots  \\
z & \equiv y_n \mod x_n.
\end{align*}
\end{thm}

Kita bakal nggunakake Teorema Sunzi kaya mangkene: yen $x_0$,
\dots,~$x_n$ luwih gedhe tinimbang $y_0$, \dots,~$y_n$ miturut
urutane, $z$ bisa dadi kode urutan $\tuple{y_0, \dots,y_n}$. Kanggo
nemokake maneh~$y_i$, cukup para $z$ karo~$x_i$ banjur jupuk
sisane. Supaya bisa nggunakake pangodean iki, kita kudu nemokake
biji kang cocog kanggo $x_0$, \dots,~$x_n$.

Sawetara pangamatan bakal mbantu kita. Yen diwenehi
$y_0$, \dots,~$y_n$, tetepna
\begin{align*}
j &= \max(n, y_0 + 1, \dots, y_n + 1), \\
m &= \lcm(1,\dots,j),
\end{align*}
lan tetepna
\begin{align*}
x_0 & = 1 + m \\
x_1 & = 1 + 2 \cdot m \\
x_2 & = 1 + 3 \cdot m \\
& \vdots  \\
x_n & = 1 + (n+1) \cdot m
\end{align*}
Mula rong pratelan iki bener:
\begin{enumerate}
\item\ollabel{rel-prime} $x_0,\dots,x_n$ relatif prima saben pasangan.
\item\ollabel{less} Kanggo saben $i$, $y_i < x_i$.
\end{enumerate}
Kanggo ndeleng yen \olref{rel-prime} bener, jupuk $i\ne k$ lan gatekna yen $p$ wilangan prima
lan $p \mid x_i$ sarta $p \mid x_k$, mula $p \mid 1 + (i+1) m$ lan
$p \mid 1 + (k+1) m$. Nanging banjur $p$ mbagi selisihe,
\[
(1 + (i+1)m) - (1+ (k+1)m) = (i-k) m.
\]
Amarga $p$ mbagi $1 + (i+1)m$, wilangan iku ora bisa uga mbagi $m$
(yen iya, pambagian kang kapisan bakal menehi sisa~$1$). Mula $p$
mbagi $i-k$, amarga $p$ mbagi $(i-k)m$. Nanging $\left|i-k\right|$ paling
gedhe~$n$, lan kita milih $j \geq n$, mula iki njalari
$p \mid m$, kang maneh dadi kontradiksi. Dadi ora ana wilangan prima kang
mbagi loro-lorone $x_i$ lan $x_k$. Klausa~\olref{less} gampang:
kita nduweni $y_i < j \leq m < x_i$.

Saiki ayo mbuktekake lema fungsi $\beta$. Elinga yen kita bisa nggunakake
$0$, fungsi penerus, panjumlahan, perkalian, $\Char{=}$, proyeksi, lan
saben fungsi kang ditegesi saka fungsi mau liwat komposisi lan
panggolekan tanpa wates kang ditrapake marang fungsi reguler. Kita uga bisa
nggunakake relasi yen fungsi karakteristike bisa ditegesi kanthi
cara mangkono. Kaya sadurunge, kita bisa nuduhake yen relasi-relasi
iki katutup tumrap kombinasi Boolean lan kuantifikasi winates;
umpamane:
\begin{align*}
\fn{not}(x) & \defis \Char{=}(x,0)\\
\bmin{x \leq z}{R(x,y)} & \defis \umin{x}{(R(x,y) \lor x = z)}\\
\bexists{x \leq z}{R(x,y)} & \defiff R(\bmin{x \leq z}{R(x,y)}, y)
\end{align*}
Banjur kita bisa nuduhake yen kabeh iki uga bisa ditegesi tanpa
rekursi primitif:
\begin{enumerate}
\item Fungsi pasangan, $J(x,y) = \frac{1}{2}[(x+y)(x+y+1)] + x$;
% maybe explain more what is going on here, a bit confusing.
\item fungsi proyeksi
\begin{align*}
K(z) & = \bmin{x \leq z}{\bexists{y \leq z}{z = J(x,y)}},\\
L(z) & = \bmin{y \leq z}{\bexists{x \leq z}{z = J(x,y)}};
\end{align*}
\item relasi luwih cilik tinimbang $x < y$;
\item relasi bisa dipara $x \mid y$;
% \item $x \tsub y$
% \item $\fn{Prime}(x)$
% \item Assuming $p$ is prime, the relation ``$x$ is a power of $p$'':
% \[
% \bforall{y \leq x}{(y \mid x \lif y = 1 \lor y = x)}.
% \]
\item fungsi $\fn{rem}(x,y)$ kang mbalekake sisa nalika
  $y$ dipara karo~$x$.
\end{enumerate}
Cathetan panjelasan SOL6-F085: nalika divisor nol, rem bisa diwenehi asil nol minangka panjangkepan total; konstruksi beta mung nggunakake divisor positif. Modulus positif ing definisi congruence lan pasangan indeks beda ing bukti relatif prima saiki diterangake kanthi eksplisit.
Saiki tegesna
\begin{align*}
\beta^*(d_0,d_1,i) & = \fn{rem}(1+(i+1) d_1,d_0) \text{ lan}\\
\beta(d,i) & = \beta^*(K(d),L(d),i).
\end{align*}
Rumus kapisan netepake beta lintang, lan rumus kapindho netepake beta. Iki fungsi kang kita karepake. Yen diwenehi $a_0,\dots,a_n$ kaya ing ndhuwur, tetepna
\[
j = \max(n,a_0+1,\dots,a_n+1),
\]
lan tetepna $d_1 = \lcm(1,\dots,j)$. Miturut \olref{rel-prime} ing ndhuwur,
kita ngerti yen $1+d_1$, $1+2 d_1$, \dots, $1+(n+1) d_1$ relatif prima
saben pasangan, lan miturut~\olref{less} kabeh luwih gedhe tinimbang
$a_0,\dots,a_n$ miturut urutane. Miturut Teorema Sunzi ana biji~$d_0$
saengga kanggo saben~$i$,
\[
d_0 \equiv a_i \mod (1+(i+1)d_1)
\]
mula (amarga $d_1$ luwih gedhe tinimbang~$a_i$),
\[
a_i = \fn{rem}(1+(i+1)d_1,d_0).
\]
Tetepna $d = J(d_0,d_1)$. Mula kanggo saben $i \le n$, kita nduweni
\begin{align*}
\beta(d,i) & =  \beta^*(d_0,d_1,i) \\
& =  \fn{rem}(1+(i+1) d_1,d_0) \\
& =  a_i
\end{align*}
kaya kang dibutuhake. Iki ngrampungake bukti lema
fungsi~$\beta$.

\begin{prob}
  Tuduhna yen relasi $x < y$, $x \mid y$, lan
  fungsi~$\fn{rem}(x,y)$ bisa ditegesi tanpa rekursi primitif.
  Kowé oleh nggunakake $0$, fungsi penerus, panjumlahan, perkalian,
  $\Char{=}$, proyeksi, lan panggolekan lan kuantifikasi winates.
\end{prob}

\end{document}
